% GENERATED FILE. Edit canonical lessons, then run npm run generate:books.
% canonical-source-hash: fnv1a64:2a274e5c91cbbbae
% canonical-lessons: SA-C512-churika, SA-C512-krntati, SA-C512-curnayati, SA-C512-surpam, SA-C512-samyojayati

\chapter{Chop and Grind}
\label{ch:sa-chop-and-grind}
\hlchaptermodality{512}

\begin{chapteropening}
\textbf{By the end of this chapter:} \emph{I can describe getting food ready to cook.}

Complete the last lesson of chapter 512: I can describe getting food ready to cook.
\end{chapteropening}

\section[\texorpdfstring{\sk{छुरिका} (churikā)}{छुरिका (churikā)}]{\sk{छुरिका} (churikā) — a knife}

\label{lesson:SA-C512-churika}



\begin{quote}
Before the new one: say the Sanskrit for cooking, then the Sanskrit for roasts.
\end{quote}

\subsection*{You'll want to know: \sk{छुरिका}}
\textbf{\sk{छुरिका}} — \emph{churikā} — \textquotedblleft{}a knife\textquotedblright{}.

\begin{quote}
\sk{माता} \sk{छुरिकया} \sk{शाकं} \sk{छिनत्ति।}
\end{quote}

\begin{quote}
\emph{mātā churikayā śākaṃ chinatti.}
\end{quote}

\begin{quote}
Mother cuts the vegetables with a knife.
\end{quote}

A feminine noun.

\begin{cousinweb}
\textbf{\sk{छुरिका}} (\emph{churikā}) comes from \emph{kṣura}, \textquotedblleft{}a razor\textquotedblright{}, as MW derives it; MW records \emph{churikā} as a knife (MW).
\end{cousinweb}

\begin{grammarlens}[title={What you've built}]
One more word for describing what you lived and felt.
\end{grammarlens}

\subsection*{Your turn}
\begin{itemize}
\raggedright
  \item \emph{Say it:} \emph{churikā}
  \item \emph{Say it:} \emph{churikā}, once more
  \item \emph{Say it:} \emph{mātā churikayā śākaṃ chinatti.} — Mother cuts the vegetables with a knife.
  \item \emph{Recall:} say the Sanskrit for cooking, then the Sanskrit for roasts, then say \emph{churikā} again
\end{itemize}

\begin{tcolorbox}[breakable,colback=teal!4,colframe=teal!35!black,title={Before you move on}]
What does \sk{छुरिका} mean? (\textquotedblleft{}A knife\textquotedblright{}.) Say it once more.
\end{tcolorbox}
\section[\texorpdfstring{\sk{कृन्तति} (kṛntati)}{कृन्तति (kṛntati)}]{\sk{कृन्तति} (kṛntati) — cuts up}

\label{lesson:SA-C512-krntati}



\begin{quote}
Before the new one: say the Sanskrit for roasts, then the Sanskrit for a knife.
\end{quote}

\subsection*{You'll want to know: \sk{कृन्तति}}
\textbf{\sk{कृन्तति}} — \emph{kṛntati} — \textquotedblleft{}cuts up\textquotedblright{}.

\begin{quote}
\sk{सः} \sk{आलुकानि} \sk{कृन्तति।}
\end{quote}

\begin{quote}
\emph{saḥ ālukāni kṛntati.}
\end{quote}

\begin{quote}
He cuts up the potatoes.
\end{quote}

Past (he …): \sk{सः} \sk{पलाण्डुं} \sk{कृत्तवान्}, \emph{saḥ palāṇḍuṃ kṛttavān}.

\begin{cousinweb}
\textbf{\sk{कृन्तति}} (\emph{kṛntati}) comes from the root √kṛt, \textquotedblleft{}to cut, to cut in pieces, to divide\textquotedblright{}, whose present is \emph{kṛntati}; its past participle is \emph{kṛtta} (MW).
\end{cousinweb}

\begin{grammarlens}[title={What you've built}]
One more word for describing what you lived and felt.
\end{grammarlens}

\subsection*{Your turn}
\begin{itemize}
\raggedright
  \item \emph{Say it:} \emph{kṛntati}
  \item \emph{Say it:} \emph{kṛntati}, once more
  \item \emph{Say it:} \emph{saḥ ālukāni kṛntati.} — He cuts up the potatoes.
  \item \emph{Recall:} say the Sanskrit for roasts, then the Sanskrit for a knife, then say \emph{kṛntati} again
\end{itemize}

\begin{tcolorbox}[breakable,colback=teal!4,colframe=teal!35!black,title={Before you move on}]
What does \sk{कृन्तति} mean? (\textquotedblleft{}Cuts up\textquotedblright{}.) Say it once more.
\end{tcolorbox}
\section[\texorpdfstring{\sk{चूर्णयति} (cūrṇayati)}{चूर्णयति (cūrṇayati)}]{\sk{चूर्णयति} (cūrṇayati) — grinds to powder}

\label{lesson:SA-C512-curnayati}



\begin{quote}
Before the new one: say the Sanskrit for a knife, then the Sanskrit for cuts up.
\end{quote}

\subsection*{You'll want to know: \sk{चूर्णयति}}
\textbf{\sk{चूर्णयति}} — \emph{cūrṇayati} — \textquotedblleft{}grinds to powder\textquotedblright{}.

\begin{quote}
\sk{सा} \sk{मरीचं} \sk{चूर्णयति।}
\end{quote}

\begin{quote}
\emph{sā marīcaṃ cūrṇayati.}
\end{quote}

\begin{quote}
She grinds the pepper.
\end{quote}

Past (she …): \sk{सा} \sk{मरीचं} \sk{चूर्णितवती}, \emph{sā marīcaṃ cūrṇitavatī}.

\begin{cousinweb}
\textbf{\sk{चूर्णयति}} (\emph{cūrṇayati}) is built on \emph{cūrṇa}, \textquotedblleft{}powder, flour\textquotedblright{}: MW gives \emph{cūrṇayati} as to reduce to powder, to grind, to pound (MW).
\end{cousinweb}

\begin{grammarlens}[title={What you've built}]
One more word for describing what you lived and felt.
\end{grammarlens}

\subsection*{Your turn}
\begin{itemize}
\raggedright
  \item \emph{Say it:} \emph{cūrṇayati}
  \item \emph{Say it:} \emph{cūrṇayati}, once more
  \item \emph{Say it:} \emph{sā marīcaṃ cūrṇayati.} — She grinds the pepper.
  \item \emph{Recall:} say the Sanskrit for a knife, then the Sanskrit for cuts up, then say \emph{cūrṇayati} again
\end{itemize}

\begin{tcolorbox}[breakable,colback=teal!4,colframe=teal!35!black,title={Before you move on}]
What does \sk{चूर्णयति} mean? (\textquotedblleft{}Grinds to powder\textquotedblright{}.) Say it once more.
\end{tcolorbox}
\section[\texorpdfstring{\sk{शूर्पम्} (śūrpam)}{शूर्पम् (śūrpam)}]{\sk{शूर्पम्} (śūrpam) — a winnowing basket}

\label{lesson:SA-C512-surpam}



\begin{quote}
Before the new one: say the Sanskrit for cuts up, then the Sanskrit for grinds to powder.
\end{quote}

\subsection*{You'll want to know: \sk{शूर्पम्}}
\textbf{\sk{शूर्पम्}} — \emph{śūrpam} — \textquotedblleft{}a winnowing basket\textquotedblright{}.

\begin{quote}
\sk{माता} \sk{शूर्पेण} \sk{तण्डुलान्} \sk{शोधयति।}
\end{quote}

\begin{quote}
\emph{mātā śūrpeṇa taṇḍulān śodhayati.}
\end{quote}

\begin{quote}
Mother cleans the rice with a winnowing basket.
\end{quote}

A neuter noun.

\begin{cousinweb}
\textbf{\sk{शूर्पम्}} (\emph{śūrpam}) is the word MW records for a winnowing basket, the wicker tray that is shaken to fan the husk away; MW gives no derivation, so its origin is uncertain (MW).
\end{cousinweb}

\begin{grammarlens}[title={What you've built}]
One more word for describing what you lived and felt.
\end{grammarlens}

\subsection*{Your turn}
\begin{itemize}
\raggedright
  \item \emph{Say it:} \emph{śūrpam}
  \item \emph{Say it:} \emph{śūrpam}, once more
  \item \emph{Say it:} \emph{mātā śūrpeṇa taṇḍulān śodhayati.} — Mother cleans the rice with a winnowing basket.
  \item \emph{Recall:} say the Sanskrit for cuts up, then the Sanskrit for grinds to powder, then say \emph{śūrpam} again
\end{itemize}

\begin{tcolorbox}[breakable,colback=teal!4,colframe=teal!35!black,title={Before you move on}]
What does \sk{शूर्पम्} mean? (\textquotedblleft{}A winnowing basket\textquotedblright{}.) Say it once more.
\end{tcolorbox}
\section[\texorpdfstring{\sk{संयोजयति} (saṃyojayati)}{संयोजयति (saṃyojayati)}]{\sk{संयोजयति} (saṃyojayati) — combines}

\label{lesson:SA-C512-samyojayati}



\begin{quote}
Before the new one: say the Sanskrit for grinds to powder, then the Sanskrit for a winnowing basket.
\end{quote}

\subsection*{You'll want to know: \sk{संयोजयति}}
\textbf{\sk{संयोजयति}} — \emph{saṃyojayati} — \textquotedblleft{}combines\textquotedblright{}.

\begin{quote}
\sk{सा} \sk{दधि} \sk{शर्करया} \sk{संयोजयति।}
\end{quote}

\begin{quote}
\emph{sā dadhi śarkarayā saṃyojayati.}
\end{quote}

\begin{quote}
She combines the curd with sugar.
\end{quote}

Past (he …): \sk{सः} \sk{शर्करां} \sk{दुग्धेन} \sk{संयोजितवान्}, \emph{saḥ śarkarāṃ dugdhena saṃyojitavān}.

\begin{cousinweb}
\textbf{\sk{संयोजयति}} (\emph{saṃyojayati}) is built on \emph{sam}, \textquotedblleft{}together\textquotedblright{}, and the causative of the root √yuj, \textquotedblleft{}to join\textquotedblright{}, which also gives \emph{yojayati}; MW gives \emph{saṃyuj} as to join together, to combine (MW).
\end{cousinweb}

\begin{grammarlens}[title={What you've built}]
One more word for describing what you lived and felt.
\end{grammarlens}

\subsection*{Your turn}
\begin{itemize}
\raggedright
  \item \emph{Say it:} \emph{saṃyojayati}
  \item \emph{Say it:} \emph{saṃyojayati}, once more
  \item \emph{Say it:} \emph{sā dadhi śarkarayā saṃyojayati.} — She combines the curd with sugar.
  \item \emph{Recall:} say the Sanskrit for grinds to powder, then the Sanskrit for a winnowing basket, then say \emph{saṃyojayati} again
\end{itemize}

\begin{tcolorbox}[breakable,colback=teal!4,colframe=teal!35!black,title={Before you move on}]
What does \sk{संयोजयति} mean? (\textquotedblleft{}Combines\textquotedblright{}.) Say it once more.
\end{tcolorbox}
